\subsection{An inseparable multisection construction is also excluded}
\begin{proposition}[A2P: the relative Frobenius multisection obstruction]
Let $S$ be a smooth projective integral variety of positive dimension
over an algebraically closed field of characteristic $p>0$.
Use the quotient convention for $P=\mathbf P_S(E)$, where $E$ has
rank two. Let $F_S^*E\twoheadrightarrow N$ be a line quotient.
The relative Frobenius
$P\to\mathbf P_S(F_S^*E)$ has degree $p$ on each projective-line
fibre. Let $Y\subset P$ be the inverse image of the section
defined by that quotient, and suppose $Y$ is smooth and integral.
If $Y$ is ample, then
\[
       \Pic(S)[p](k)\longrightarrow\Pic(Y)[p](k)
\]
is injective. Thus this construction cannot give an A2 example
by normalizing a nontrivial $\mu_p$-torsor on $S$.
\end{proposition}
\begin{proof}
Write $\rho:Y\to S$ and $K=\OO_P(1)|_Y$. The morphism $\rho$
is finite flat and purely inseparable of degree $p$. The section
of $\mathbf P_S(F_S^*E)$ is cut out by its kernel line bundle
$(\det E)^p\otimes N^{-1}$. Pulling back its divisor gives
\[
 \OO_P(Y)=\OO_P(p)\otimes
                 \pi^*((\det E)^{-p}\otimes N),
 \qquad K^p=\rho^*N.                                  \tag{A2P.1}
\]
Put $Q=K^2\otimes\rho^*(\det E)^{-1}$. The projective-bundle
canonical formula and smooth-divisor adjunction give
\[
 \OO_P(Y)|_Y=Q^p,\qquad
 \omega_Y\otimes(\rho^*\omega_S)^{-1}=Q^{p-1}.          \tag{A2P.2}
\]
Indeed $\omega_P=\OO_P(-2)\otimes\pi^*(\omega_S\otimes\det E)$;
substituting (A2P.1) gives exponents $2p-2$ on $K$ and $1-p$
on $\rho^*\det E$, exactly as in (A2P.2).

Suppose a nontrivial $M\in\Pic(S)[p](k)$ becomes trivial on $Y$.
As in A2Z, its integral $\mu_p$-torsor $T\to S$ has degree $p$.
A compatible trivialization of $\rho^*M$ gives a finite birational
map $Y\to T$, so $Y$ is its normalization. The Gorenstein torsor
has $\omega_T\simeq\omega_S|_T$, and finite duality gives
\[
        \omega_Y\otimes(\rho^*\omega_S)^{-1}
                    \simeq\OO_Y(-D)                         \tag{A2P.3}
\]
for an effective conductor divisor $D$. The precise conductor
argument, including nonnormal $T$, was proved in A2Z using
\cite[Lemma 47.15.8]{DGJ_StacksFiniteDuality}.
Equations (A2P.2)--(A2P.3) imply $Q^{p-1}\simeq\OO_Y(-D)$.
If $Y$ is ample in $P$, its normal line $Q^p$ is ample, hence
$Q$ is ample. An ample positive power cannot be the negative of
an effective divisor on the positive-dimensional projective
variety $Y$: intersection with any ample class of complementary
degree is positive on the former and nonpositive on the latter.
This contradiction proves injectivity.
\end{proof}

This is an exclusion for the specified relative Frobenius
multisection. It neither assumes a disjoint infinity section nor
classifies arbitrary smooth ample divisors in projective bundles.
The unrestricted A2 question remains undecided.
